\documentclass{article}
\title{Deriving the Oseen-Lamb Vortex: Three Methods}
\author{T.M. Montgomery}
\date{July 13, 2025 \\ Revised: July 20, 2025}

\usepackage{amsmath}
\usepackage{graphicx}
\usepackage[colorlinks=true, allcolors=blue]{hyperref}
\usepackage[letterpaper,top=2cm,bottom=2cm,left=3cm,right=3cm,marginparwidth=1.75cm]{geometry}
\usepackage{amsthm}
\usepackage{amssymb}
\usepackage{graphicx}
\usepackage{subcaption}
\usepackage[export]{adjustbox}
\usepackage{wrapfig}
\usepackage{xcolor}

\usepackage[nottoc]{tocbibind}

\begin{document}
\maketitle

\tableofcontents

\begin{abstract}
    In great depth, we investigate three methods to derive the axially symmetric azimuthal velocity model: the Oseen-Lamb vortex. The first method is finding a velocity solution of the form, $u_\theta =\frac{\Gamma_0}{2\pi r}g(r,t)$ for some function, $g:\mathbb{R}^2\rightarrow\mathbb{R}^2$. The second method is to solve the vorticity transport equation in polar coordinates, which yields a Sturm-Liouville differential equation with a lengthy process. The third method is by no means as robust as the first two, but it utilizes Green's theorem and a velocity parcel to obtain a separable DE in terms of circulation.
\end{abstract}

\section*{Introduction \& Boundary Conditions:}
The solution to the advectionless and pressureless Navier-Stokes equation given a unidirectional velocity field, $\vec{u}=\left< 0,u_\theta ,0\right>$, is the Lamb-Oseen vortex.
\begin{align}\label{Eq1}
u_{\theta} (r,t) = \frac{\Gamma_0}{2\pi r}\left(1-e^{-\frac{r^2}{4\nu t}}\right) \:\:\:\:,\:\: \Gamma_0 \:,\: \nu \in \mathbb{R}^+
\end{align}
with the boundary conditions, 
\begin{align*}
    \left\{\begin{matrix}
    u(0,t)= 0 \:\:,& u(r\to \infty ,t)=0 \\ u(r,0)=\frac{\Gamma_0}{2\pi r} \:\:,& \:\: u(r,t\to \infty)=0
    \end{matrix}\right\}
\end{align*}

\section{The First Method: Finding $g(r,t)$ And Using A Similarity Variable}
Consider an azimuthal velocity solution of the form, 
\begin{align}\label{Eq2}
    u_\theta (r,t)=\frac{\Gamma_0}{2\pi r}g(r,t)
\end{align}
and substitute (\ref{Eq2}) into the laminarized momentum equation,
\begin{align}\label{Eq3}
\begin{split}
    \frac{\partial u_\theta}{\partial t}&=\nu \left(\frac{\partial^2 u_\theta}{\partial r^2}+\frac{1}{r}\frac{\partial u_\theta}{\partial r} -\frac{u_\theta}{r^2}\right)\\
    \frac{\partial}{\partial t}\left( \frac{\Gamma_0}{2\pi r}g(r,t)\right) &=\nu \left(\frac{\partial^2}{\partial r^2}\left( \frac{\Gamma_0}{2\pi r}g(r,t)\right)+\frac{1}{r}\frac{\partial}{\partial r}\left( \frac{\Gamma_0}{2\pi r}g(r,t)\right) -\frac{1}{r^2}\frac{\Gamma_0}{2\pi r}g(r,t)\right)\\
    \frac{\Gamma_0}{2\pi r}\frac{\partial g}{\partial t} &=\nu \left(\left( 2\frac{\Gamma_0}{2\pi r^3}g -\frac{\Gamma_0}{2\pi r^2}\frac{\partial g}{\partial r} -\frac{\Gamma_0}{2\pi r^2}\frac{\partial g}{\partial r}+\frac{\Gamma_0}{2\pi r}\frac{\partial^2 g}{\partial r^2}  \right)+\frac{1}{r} \left( -\frac{\Gamma_0}{2\pi r^2}g +\frac{\Gamma_0}{2\pi r}\frac{\partial g}{\partial r} \right) -\frac{1}{r^2}\frac{\Gamma_0}{2\pi r}g\right)\\
    \frac{\partial g}{\partial t} &=\nu \left(\left( \frac{2}{r^2}g -\frac{2}{r}\frac{\partial g}{\partial r} + \frac{\partial^2 g}{\partial r^2}  \right)+\frac{1}{r} \left( -\frac{1}{r}g +\frac{\partial g}{\partial r} \right) -\frac{1}{r^2}g\right)\\
    \frac{\partial g}{\partial t} &=\nu \left( \frac{\partial^2 g}{\partial r^2} - \frac{1}{r}\frac{\partial g}{\partial r} \right)
\end{split}
\end{align}
We can invoke a similarity solution, $g(r,t)=f(\eta)$ with $\eta =\frac{r}{\sqrt{\nu t}}$, and substitute to obtain the following derivatives:
\begin{align}\label{Eq4}
\begin{split}
    \frac{\partial g}{\partial t} &= \frac{\partial f}{\partial \eta}\frac{\partial \eta}{\partial t} = -f' \frac{r t^{-3/2}}{2\sqrt{\nu}}= -f' \frac{\eta}{2t}\\
    \frac{\partial g}{\partial r}&=\frac{\partial f}{\partial \eta}\frac{\partial \eta}{\partial r} = f' \frac{1}{\sqrt{\nu t}}\\
    \frac{\partial^2 g}{\partial r^2} &=\frac{\partial }{\partial r}\left( \frac{\partial f}{\partial \eta}\frac{\partial \eta}{\partial r} \right) \\
    &=\frac{\partial f}{\partial \eta}\underbrace{\frac{\partial}{\partial r}\left ( \frac{\partial \eta}{\partial r} \right ) }_{=0}
    + \frac{\partial \eta}{\partial r}\frac{\partial }{\partial r}\left ( \frac{\partial f}{\partial \eta} \right ) \\
    &= \frac{\partial \eta}{\partial r} \cdot \frac{\partial^2 f}{\partial \eta^2} \frac{\partial \eta}{\partial r}\\
    &= f'' \frac{1}{\nu t}
\end{split}
\end{align}
The following separable ordinary differential equation is achieved by the substitution of the derivatives in (\ref{Eq4}) into (\ref{Eq3}),
\begin{align*}
    -f'\frac{\eta}{2t}=\nu \left( \frac{f''}{\nu t}-\frac{f'}{r\sqrt{\nu t}} \right)\\
    f'' + \left( \frac{\eta}{2}-\frac{1}{\eta} \right)f' =0
\end{align*}
with the solution,
\begin{align*}
    \frac{d^2 f}{d\eta^2} &=-\left( \frac{\eta}{2}-\frac{1}{\eta} \right)\frac{df}{d\eta} \:\:\:,\:\:\: u(\eta)=f'(\eta)\\
    \frac{d u}{d\eta} &=-\left( \frac{\eta}{2}-\frac{1}{\eta} \right)u \\
    \int \frac{d u}{u} &= -\int \left( \frac{\eta}{2}-\frac{1}{\eta} \right)d\eta \\
    \ln(u) &= - \left( \frac{\eta^2}{4}-\ln(\eta) \right)+C \\
    u(\eta) &= C_1e^{-\frac{\eta^2}{4}+\ln(\eta)} \\
    f'(\eta) &= C_1\eta \: e^{-\frac{\eta^2}{4}} \\
    f(\eta) &= C_1 \int \eta\: e^{-\frac{\eta^2}{4}} d\eta \:\:\:,\:\:\: p=-\frac{\eta^2}{4}\:\:\:,\:\;\: dp=-\frac{\eta}{2}d\eta \\
    f(\eta) &= C_1 \int \eta\: e^{p} \left( -\frac{2}{\eta} \right) dp \\
    f(\eta) &= A e^{-\frac{\eta^2}{4}} +B \:\:\:,\:\:\: A =-2C_1\\
    \therefore \: \omega_{\hat{z}}(r,t) &= A e^{-\frac{r^2}{4\nu t}}+B
\end{align*}
The constants, $A$ and $B$, can be found by the circulation formula, defined as the vorticity flux integral over a circular area differential.
\begin{align*}
    \Gamma = \int_A \vec{\omega}_{\hat{z}}\cdot dA &= \int_{0}^{\infty} \omega_{\hat{z}} \cdot 2\pi rdr \\
    &= 2\pi\int_{0}^{\infty} \left( A e^{-\frac{r^2}{4\nu t}}+B \right) r \: dr \\
    &= 2\pi \left( \int_{0}^{\infty} Ar e^{-\frac{r^2}{4\nu t}} \: dr +\int_{0}^{\infty} Br\:dr\right) \\
    &= 2\pi \left( \int_{0}^{\infty} Ar \left( \frac{2\nu t}{r}\right) e^{-u} \: du + \left. B\frac{r^2}{2}\right|_{0}^{\infty} \right) \\
    &= 4\pi \nu A t\int_{0}^{\infty} e^{-u} \: du + \pi \left. B r^2\right|_{0}^{\infty}  \\
    &= 4\pi \nu A t + \pi \left. B r^2\right|_{0}^{\infty} \:\:\:,\:\:\: B=0  \\
    \Gamma_0 &= 4\pi \nu A t \\
    \therefore A &= \frac{\Gamma_0}{4\pi \nu t}
\end{align*}
where $\Gamma_0 \in \mathbb{R}^+$ is the initial circulation and $B=0$ since there are no regions in the fluid with infinite circulation. The vorticity becomes the Gaussian distribution,
\begin{align}\label{Eq5}
    \omega_{\hat{z}}(r,t) &= \frac{\Gamma_0}{4\pi \nu t} e^{-\frac{r^2}{4\nu t}}
\end{align}
Because $\omega_{\hat{z}} = \nabla \times u_\theta =\left(\frac{\partial u_\theta}{\partial r}+\frac{u_\theta}{r}\right)\hat{z}$, we can use the integrating factor, $e^{\int \frac{1}{r} dr}=r$, to isolate $u_\theta$ and substitute Equation (\ref{Eq5}) as follows:
\begin{align}\label{Eq6}
\begin{split}
    u_\theta (r,t) &= \frac{1}{r}\int r\omega_{\hat{z}}\: dr \\
    &= \frac{\Gamma_0}{4\pi \nu t} \frac{1}{r}\int r e^{-\frac{r^2}{4\nu t}} dr \\
    &= \frac{\Gamma_0}{4\pi \nu t} \frac{1}{r}\int r e^{u} \left(-\frac{2\nu t}{r} \right) du \\
    &= -\frac{\Gamma_0}{2\pi} \frac{1}{r}\int e^{u} du \\
    &= -\frac{\Gamma_0}{2\pi r} e^{-\frac{r^2}{4\nu t}} +C \\
\end{split}
\end{align}
Using the boundary condition, $u_{\theta}(0,t)=0$, $C=\frac{\Gamma_0}{2\pi r}$, the solution to Navier-Stokes is obtained.
\begin{align*}
    u_\theta (r,t) = \frac{\Gamma_0}{2\pi r}\left( 1-e^{-\frac{r^2}{4\nu t}}\right) 
\end{align*}

\section{The Second Method: Using The Vorticity Transport Equation And The Frobenius Method}
Taking the curl of the momentum equation,
\begin{align*}
    \frac{\partial u_\theta}{\partial t}&=\nu \left(\frac{\partial^2 u_\theta}{\partial r^2}+\frac{1}{r}\frac{\partial u_\theta}{\partial r} -\frac{u_\theta}{r^2}\right)
\end{align*}
we have the $\hat{z}$ directional vorticity transport equation in $\mathbb{R}^2$:
\begin{align}\label{Eq7}
    \boldsymbol{\hat{z}}\:\: : \:\: \frac{\partial \omega}{\partial t}=\nu \left(\frac{\partial^2 \omega}{\partial r^2}+\frac{1}{r}\frac{\partial \omega}{\partial r}\right)
\end{align}
We invoke a new similarity solution, $\eta =\frac{r}{\sqrt{4\nu t}}$ for $\omega_{\hat{z}} (r,t)=\frac{1}{t}f(\eta)$, and find each derivative:
\begin{align}\label{Eq8}
\begin{split}
    \frac{\partial \omega}{\partial t} &= \frac{\partial }{\partial t}\left(\frac{f}{t} \right)=\frac{\partial}{\partial t}\left(\frac{1}{t}\right) +\frac{\partial f}{\partial \eta}\frac{\partial \eta}{\partial t} = -\frac{1}{t^2}f-\frac{2\nu r}{t(4\nu t)^{3/2}}f'\\
    \frac{\partial \omega}{\partial r} &=\frac{1}{t}\frac{\partial f}{\partial \eta}\frac{\partial \eta}{\partial r} = f' \frac{1}{t\sqrt{4\nu t}}\\
    \frac{\partial^2 \omega}{\partial r^2} &=\frac{\partial }{\partial r}\left( \frac{1}{t}\frac{\partial f}{\partial \eta}\frac{\partial \eta}{\partial r} \right) \\
    &=\frac{1}{t}\frac{\partial f}{\partial \eta}\underbrace{\frac{\partial}{\partial r}\left ( \frac{\partial \eta}{\partial r} \right ) }_{=0}
    + \frac{1}{t}\frac{\partial \eta}{\partial r}\frac{\partial }{\partial r}\left ( \frac{\partial f}{\partial \eta} \right ) \\
    &= \frac{1}{t}\frac{\partial \eta}{\partial r} \cdot \frac{\partial^2 f}{\partial \eta^2} \frac{\partial \eta}{\partial r}\\
    &= f'' \frac{1}{4\nu t^2}
\end{split}
\end{align}
By substitution into (\ref{Eq7}), we have the ODE:
\begin{align*}
    -\frac{1}{t^2}f-\frac{2\nu r}{t(4\nu t)^{3/2}}f' &= \nu \left(f'' \frac{1}{4\nu t^2}+ f' \frac{1}{r}\frac{1}{t\sqrt{4\nu t}} \right)\\
    f'' + \left(\frac{1}{\eta}+2\eta \right)f' +4f &= 0
\end{align*}
By inspection, we have a homogeneous Sturm-Liouville DE \cite{13.2Sturm-Liouville} of the form,
\begin{align*}
    \frac{d}{dx}\left[ p(x)\frac{d\phi}{dx}\right]+\left[q(x)+\lambda w(x) \right]\phi(x)=0
\end{align*}
with a general solution in the form of a linearly independent eigenfunction pair of solutions to the DE,
\begin{align*}
    \phi (x)=c_1\phi_1 (x)+c_2 \phi_2(x)
\end{align*}
To demonstrate with the integrating factor, $p(\eta)=e^{\int \frac{1}{\eta}+2\eta \:d\eta} = \eta e^{\eta^2}$, we can get our DE in Sturm-Liouville form (where $q(x)=0$),
\begin{align*}
    \eta e^{\eta^2}f'' + \left(\frac{1}{\eta}+2\eta \right)\eta e^{\eta^2}f' +4\eta e^{\eta^2}f &= 0 \\
    \frac{d}{d\eta}\left( \eta e^{\eta^2}f'\right) +4\eta e^{\eta^2}f &= 0
\end{align*}
and a general solution, $f(\eta)=c_1 f_1(\eta)+c_2 f_2(\eta)$. Using Frobenius's method \cite{4.4.2FrobeniusMethod}, let our DE be,
\begin{align*}
    y'' + \left(\frac{1}{x}+2x \right)y' +4y = 0
\end{align*}
with the solution in the form, $y(x)=c_1 y_1 (x)+c_2 y_2 (x)$. Each derivative can be expressed as the series,
\begin{align*}
    y(x) &= \sum_{k=0}^{\infty}c_k x^{k+r}\\
    y'(x) &= \sum_{k=0}^{\infty}c_k (k+r)x^{k+r-1}\\
    y''(x) &= \sum_{k=0}^{\infty}c_k (k+r)(k+r-1)x^{k+r-2}
\end{align*}
By substitution,
\begin{align*}
    0&= \sum_{k=0}^{\infty}c_k (k+r)(k+r-1)x^{k+r-2} + \left( \frac{1}{x}+2x\right)\sum_{k=0}^{\infty}c_k (k+r)x^{k+r-1} + 4\sum_{k=0}^{\infty}c_k x^{k+r}\\
    &= \sum_{k=0}^{\infty}c_k (k+r)(k+r-1)x^{k+r-2} +  \left\{\frac{1}{x}\sum_{k=0}^{\infty}c_k (k+r)x^{k+r-1} + 2x\sum_{k=0}^{\infty}c_k (k+r)x^{k+r-1} \right\} + 4\sum_{k=0}^{\infty}c_k x^{k+r}\\
    &= \sum_{k=0}^{\infty}c_k (k+r)(k+r-1)x^{k+r-2} + \sum_{k=0}^{\infty}c_k (k+r)x^{k+r-2} + 2\sum_{k=0}^{\infty}c_k (k+r)x^{k+r} + 4\sum_{k=0}^{\infty}c_k x^{k+r}\\
    &= \sum_{k=0}^{\infty}c_k \left[(k+r)(k+r-1) + (k+r) \right]x^{k+r-2}
    + \sum_{k=0}^{\infty} 2c_k \left(k+r + 2\right)x^{k+r} \\
    &= \sum_{k=0}^{\infty}c_k \left(k+r\right)^2x^{k+r-2}
    + \sum_{k=0}^{\infty} 2c_k \left(k+r + 2\right)x^{k+r} \\
    &= \sum_{k=0}^{\infty}c_k \left(k+r\right)^2x^{k+r-2}
    + \sum_{k-2=0}^{\infty} 2c_{k-2} \left(k-2+r + 2\right)x^{k-2+r} \\
    &= \sum_{k=0}^{\infty}c_k \left(k+r\right)^2x^{k+r-2}
    + \sum_{k=2}^{\infty} 2c_{k-2} \left(k+r\right)x^{k+r-2} \\
    &= c_0 r^2 x^{r-2}+c_1 \left(1+r\right)^2 x^{r-1}+\sum_{k=2}^{\infty} \left[ c_k \left(k+r\right)^2 + 2c_{k-2} \left(k+r\right) \right]x^{k+r-2}
\end{align*}
With $c_1=0$, we have a double root in which $r^2 =0$ and $r_1=r_2=0$, thus the indicial equation becomes,
\begin{align*}
    c_k k^2 + 2c_{k-2} k=0\\
    \therefore c_k = -\frac{2c_{k-2}}{k} \:\:,\:\: k\geq 2
\end{align*}
Testing the first few values of $c_k$, we observe the pattern,
\begin{center}
\begin{tabular}{ c c c }
 $c_1=0$ & $c_4=-\frac{2}{4}c_2 = \frac{1}{2}c_0$ & $c_7=0$ \\ 
 $c_2 =-c_0$ & $c_5=0$ & $c_8 = -\frac{2}{8}c_6 =-\frac{1}{4}\left( -\frac{1}{3}\left(\frac{1}{2}c_0 \right)\right) $ \\  
 $c_3=0$ & $c_6=-\frac{1}{3}c_4 =-\frac{1}{3}\left(-\frac{1}{2}c_2 \right)=-\frac{1}{3}\left(\frac{1}{2}c_0 \right)$ & $c_9=0$
\end{tabular}
\end{center}
Therefore, $\left\{ -c_0 , \frac{1}{2!}c_0,-\frac{1}{3!}c_0, \frac{1}{4!}c_0, -\frac{1}{5!}c_0 ,\dots \right\}$ and the first solution, $y_1 (x)$, becomes,
\begin{align*}
    y_1 (x) &= \left | x\right |^{r_{1}} \sum_{k=0}^{\infty}c_k x^{k} \:\:\:,\:\:\: c_{2m-1}=0\:\:,m\in \mathbb{N}\\
    &= (1)\left( c_0 -\frac{c_0}{1!}x^2 + \frac{c_0}{2!}x^4 - \frac{c_0}{3!}x^6 +\dots \right)\\
    &= \sum_{k=0}^{\infty}c_0\frac{(-x^2)^k}{k!}\\
    &=e^{-x^2} \:\:\textit{where,}\:\: c_0=1
\end{align*}
The second solution, $y_2(x)$, can be found by substituting $y_2 (x)=u(x)e^{-x^2}$ into the original DE:
\begin{align*}
    0 &= y'' + \left(\frac{1}{x}+2x \right)y' +4y \\
    &= \left[ u''e^{-x^2} - 4xu'e^{-x^2} -2ue^{-x^2}+4x^2 ue^{-x^2} \right] + \left(\frac{1}{x}+2x \right)\left[ u'e^{-x^2} -2xue^{-x^2} \right] +4ue^{-x^2} \\
    &= u'' - 4xu' -2u +4x^2 u + \left(\frac{1}{x}+2x \right)\left( u' -2xu \right) +4u \\
    &= u'' - 4xu' -2u +4x^2 u + \left( \frac{1}{x}u' -2u+2xu' - 4x^2 u \right)+4u \\
    &= u'' +\frac{1}{x}u'-2xu'  \\
    u''&= -\left(\frac{1}{x}-2x \right)u' \:\:\:,\:\:\: \textit{solve with,}\:\: w(x)=u'(x)\\
    u(x)&= \frac{C_1}{2} \text{Ei}(x^2) \:\:\:,\:\:\: C_2 =C_1 /2\\
    \therefore y_2 (x)&= C_2 \text{Ei}(x^2) e^{-x^2} 
\end{align*}
The general solution is then
\begin{align*}
    y(x) &=c_1 y_1 (x) + c_2 y_2 (x)\\
    y(x)&= c_1 e^{-x^2} + c_2 \text{Ei}(x^2) e^{-x^2} \:\:\:\textit{denoting,}\:\: y(x)=f(\eta)\:\:,\:\:x=\eta\\
    \therefore f(\eta) &= c_1 e^{-\eta^2} + c_2 \operatorname{Ei}(\eta^2) e^{-\eta^2}
\end{align*}
Because $\displaystyle \lim_{x\to 0} \operatorname{Ei}(x^2)$ is not well defined, then $c_2 =0$. Substituting $\eta$ and $f(\eta)$ into $\omega_{\hat{z}} (r,t)=\frac{1}{t}f(\eta)$ yields,
\begin{align*}
    \omega_{\hat{z}} (r,t) &= c_1 e^{-\frac{r^2}{4\nu t}}
\end{align*}
which leads to Equation (\ref{Eq5}) and the method in Equation (\ref{Eq6}).

\section{The Third Method: Green's Theorem And The Velocity Parcel in $\mathbb{R}^2$}
This method is speculative, superfluous, and by no means robust in comparison to the previous two methods explored.

\begin{figure}[h]
    \centering
    \includegraphics[width=0.9\linewidth]{figs/fig1.png}
    \caption{(Left) A general curve, $C$, enclosing a region, $D$, whose perimeter differential, $ds$, has $x$ and $y$ velocity components. (Right) A parcel of this region with those components plus the differentials.}
    \label{fig:enter-label}
\end{figure}
For a velocity vector $\vec{V}=\left< u,v\right>$ along the curve $C$, we find that by Green's theorem in Cartesian coordinates, the vorticity in the region $D$ equals the flux of the vorticity with respect to the infinitesimal area.
\begin{align*}
     \Gamma &= \oint_{C} \vec{V}\cdot ds = \oint_{C}\left ( udx+vdy \right ) =\iint_{D} \left ( \frac{\partial v}{\partial x}-\frac{\partial u}{\partial y} \right )dxdy = \iint_{D} \left\| \vec{\omega}_{\hat{z}} \right\|dA = \iint_{D} \vec{\omega}_{\hat{z}} \cdot dA
\end{align*}
Taking differentials in terms of "$\delta$," we find the variation in circulation to be,
\begin{align}\label{Eq9}
\begin{split}
    \delta \Gamma &=-\left[ u \delta x + \left(v+\frac{\partial v}{\partial x}\delta x \right) \delta y \right] + \left[ v \delta y + \left(u+\frac{\partial u}{\partial y}\delta y \right) \delta x \right] \\
    &= -\left( \frac{\partial v}{\partial x}-\frac{\partial u}{\partial y} \right)\delta x \delta y \\
    &= -\vec{\omega}_{\hat{z}} \delta A \\
    \frac{ \delta \Gamma}{\delta A}&= -\vec{\omega}_{\hat{z}}
\end{split}
\end{align}
Suppose that $\Gamma (r,t)=\vec{\omega}_{\hat{z}} (r,t)a(t)$ in Equation (\ref{Eq9}) with a time-dependent area function, $a(t)=\pi R(t)^2$, and vortex radius, $R(t)$. The derivative, $\frac{d\Gamma}{dA}$, implies that $A(r)\in \Gamma (r,t)$. Then,
\begin{align*}
    \frac{ \delta \Gamma}{\delta A}&= -\frac{\Gamma (r,t)}{a(t)}\\
    \frac{ d \Gamma}{d A}&= -\frac{\Gamma (r,t)}{a(t)}\\
    \int \frac{ d \Gamma}{\Gamma }&= -\int \frac{dA}{a(t)}\\
    \Gamma (r,t)&= e^{-\frac{A(r)}{a(t)}+C}\\
    \Gamma (r,t)&=C_1 e^{-\frac{A(r)}{a(t)}} \:\:\:,\:\;\: \Gamma (0,t)=\Gamma_0=C_1\\
    \Gamma (r,t)&=\Gamma_0 e^{-\frac{A(r)}{a(t)}}
\end{align*}
Because $\Gamma (r,t)=\pi R(t)^2 \omega_{\hat{z}}(r,t)$ and $\Gamma_0=\pi R(t)^2 \omega_{\hat{z}}(t)$,
\begin{align}\label{Eq10}
\begin{split}
    \pi R(t)^2 \omega_{\hat{z}}(r,t) &= \Gamma_0 e^{-\frac{r^2}{R(t)^2}}\\
    \omega_{\hat{z}}(r,t) &=\frac{\Gamma_0}{\pi R(t)^2} e^{-\frac{r^2}{R(t)^2}}
\end{split}
\end{align}
We can use the vorticity transport equation to solve for $R(t)$, first by taking each derivative,
\begin{align*}
    \frac{\partial \omega}{\partial t} &= \frac{\Gamma_0}{\pi R^2} e^{-\frac{r^2}{R^2}}\left( \frac{2r^2}{R^3}-\frac{2}{R} \right)\frac{dR}{dt}\\
    \frac{1}{r}\frac{\partial \omega}{\partial r}&= -\frac{\Gamma_0}{\pi R^2} e^{-\frac{r^2}{R^2}}\left( \frac{2}{R^2}\right) \\
    \frac{\partial^2 \omega}{\partial r^2}&= -\frac{\Gamma_0}{\pi R^2} e^{-\frac{r^2}{R^2}}\left( \frac{2}{R^2}-\frac{4r^2}{R^4}\right)
\end{align*}
and second, by substituting into Equation (\ref{Eq7}). We find $R(t)$ by the separable DE:
\begin{align*}
    \frac{\partial \omega}{\partial t} &=\nu \left(\frac{\partial^2 \omega}{\partial r^2}+\frac{1}{r}\frac{\partial \omega}{\partial r}\right)\\
    \frac{\Gamma_0}{\pi R^2} e^{-\frac{r^2}{R^2}}\left( \frac{2r^2}{R^3}-\frac{2}{R} \right)\frac{dR}{dt} &= \nu \left[ -\frac{\Gamma_0}{\pi R^2} e^{-\frac{r^2}{R^2}}\left( \frac{2}{R^2}-\frac{4r^2}{R^4}\right) -\frac{\Gamma_0}{\pi R^2} e^{-\frac{r^2}{R^2}}\left( \frac{2}{R^2}\right) \right]\\
    \left( \frac{2r^2}{R^3}-\frac{2}{R} \right)\frac{dR}{dt} &= \nu \left[ \frac{4r^2}{R^4}-\frac{2}{R^2} -\frac{2}{R^2} \right]\\
    \left( \frac{2r^2}{R^3}-\frac{2}{R} \right)\frac{dR}{dt} &= \nu \frac{2}{R} \left( \frac{2r^2}{R^3}-\frac{2}{R} \right)\\
    \frac{dR}{dt} &= \frac{2\nu }{R}\\
    \int RdR &= 2\nu \int dt\\
    \frac{R^2}{2} &= 2\nu t +C\\
    \therefore \:\: R(t) &= \sqrt{4\nu t +R_0^2}
\end{align*}
It is not commonly found in physics journals that the initial vortex radius, $R_0$, is used, though it has interesting implications when examining confined cyclonic flow. For the purpose of this paper, $R_0=0$ so that Equation (\ref{Eq10}) becomes Equation (\ref{Eq5}).
\begin{align*}
    \omega_{\hat{z}}(r,t) &=\frac{\Gamma_0}{4\pi \nu t} e^{-\frac{r^2}{4\nu t}}
\end{align*}

\begin{figure}[h]
    \centering
    \includegraphics[width=0.6\linewidth]{figs/fig2.png}
    \caption{A curve, $C(t)$, enclosing the circular region at time, $t$, whose area and radius is $a(t)$ and $R(t)$. }
    \label{fig:enter-label}
\end{figure}

\newpage 

\bibliographystyle{siam}
\bibliography{references}
\end{document}
